\section{Conclusion}

This paper provides the first national-scale efficiency gap analysis of algorithmic redistricting, establishing empirical benchmarks for partisan fairness evaluation. Our findings demonstrate that neutral algorithmic methods substantially reduce partisan bias—62\% lower than enacted plans—while maintaining constitutional requirements and geographic compactness. The persistent $-3.2\%$ efficiency gap in algorithmic plans reflects unavoidable Democratic urban concentration and establishes a lower bound on partisan asymmetry achievable through neutral redistricting.

These results have important implications for three constituencies. For \textbf{reformers}, algorithmic redistricting offers meaningful improvements in partisan fairness without requiring constitutional amendments or abandoning traditional redistricting principles. The 62\% bias reduction represents substantial progress toward neutral governance even if perfect proportionality remains elusive. For \textbf{courts}, efficiency gap benchmarks provide quantitative standards for evaluating partisan gerrymandering claims under state constitutions. States showing enacted efficiency gaps substantially exceeding algorithmic baselines (e.g., $+7\%$ vs. $-3\%$) exhibit evidence of human manipulation beyond geographic necessity. For \textbf{scholars}, our analysis clarifies the limits of algorithmic objectivity: neutral algorithms reduce but cannot eliminate partisan effects, and the residual $-3.2\%$ bias reflects fundamental geographic constraints rather than algorithmic failures.

Future research should extend this analysis in three directions. First, comparing efficiency gaps across different algorithmic approaches (ReCom, shortest splitline, simulated annealing) would test whether our findings are method-specific or represent general properties of neutral redistricting. Second, analyzing efficiency gaps at finer geographic resolutions (census blocks rather than tracts) would increase precision and potentially reveal whether algorithmic bias can be further reduced through more granular optimization. Third, exploring multi-member district alternatives would quantify the partisan fairness gains achievable by moving beyond single-member districts entirely.

The gerrymandering crisis identified by Chief Justice Roberts in \emph{Rucho} remains acute: partisan manipulation of district boundaries undermines democratic legitimacy and entrenches political power. While \emph{Rucho} removed federal judicial oversight, state courts and legislatures retain authority to address partisan gerrymandering. Our analysis demonstrates that algorithmic redistricting provides a practical, implementable solution that substantially reduces partisan bias while respecting constitutional constraints. The question is no longer whether neutral alternatives exist—they do, and this paper quantifies their partisan fairness advantages. The question is whether political institutions will adopt them.
